% GENERATED FILE. Edit canonical lessons, then run npm run generate:books.
% canonical-source-hash: fnv1a64:bf79a1a06bd33fb8
% canonical-lessons: UR-C312-riaya, UR-C312-hukm, UR-C312-ghulam, UR-C312-baghavat, UR-C312-lalkarna

\chapter{Rebellion}
\label{ch:ur-rebellion}
\hlchaptermodality{312}

\begin{chapteropening}
\textbf{By the end of this chapter:} \emph{I can tell how people rose against a ruler.}

Complete the last lesson of chapter 312: I can tell how people rose against a ruler.
\end{chapteropening}

\section[\texorpdfstring{\ur{رعایا} (riāyā)}{رعایا (riāyā)}]{\ur{رعایا} (riāyā) — the subjects (of a ruler)}

\label{lesson:UR-C312-riaya}



\begin{quote}
Before the new one: say the Urdu for a martyr, then the Urdu for a peace settlement.
\end{quote}

\subsection*{You'll want to know: \ur{رعایا}}
\textbf{\ur{رعایا}} — \emph{riāyā} — \textquotedblleft{}the subjects (of a ruler)\textquotedblright{}.

\begin{quote}
\ur{رعایا} \ur{بادشاہ} \ur{سے} \ur{ناراض} \ur{تھی۔}
\end{quote}

\begin{quote}
\emph{riāyā bādshāh se nārāz thī.}
\end{quote}

\begin{quote}
The king's subjects were angry with him.
\end{quote}

Feminine.

\begin{cousinweb}
\textbf{\ur{رعایا}} (\emph{riāyā}) is usually traced, through Persian, to Arabic \emph{raʿāyā}, the plural of \emph{raʿiyya}, \textquotedblleft{}a flock, the people in one's care\textquotedblright{}. Urdu treats it as one feminine body: \textbf{\ur{رعایا} \ur{ناراض} \ur{تھی}}.
\end{cousinweb}

\begin{grammarlens}[title={What you've built}]
One more word for telling what happened.
\end{grammarlens}

\subsection*{Your turn}
\begin{itemize}
\raggedright
  \item \emph{Say it:} \emph{riāyā}
  \item \emph{Say it:} \emph{riāyā}, once more
  \item \emph{Say it:} \emph{riāyā bādshāh se nārāz thī.} — The king's subjects were angry with him.
  \item \emph{Recall:} say the Urdu for a martyr, then the Urdu for a peace settlement, then say \emph{riāyā} again
\end{itemize}

\begin{tcolorbox}[breakable,colback=teal!4,colframe=teal!35!black,title={Before you move on}]
What does \ur{رعایا} mean? (\textquotedblleft{}The subjects (of a ruler)\textquotedblright{}.) Say it once more.
\end{tcolorbox}
\section[\texorpdfstring{\ur{حکم} (hukm)}{حکم (hukm)}]{\ur{حکم} (hukm) — an order; a command}

\label{lesson:UR-C312-hukm}



\begin{quote}
Before the new one: say the Urdu for a peace settlement, then the Urdu for the subjects (of a ruler).
\end{quote}

\subsection*{You'll want to know: \ur{حکم}}
\textbf{\ur{حکم}} — \emph{hukm} — \textquotedblleft{}an order; a command\textquotedblright{}.

\begin{quote}
\ur{بادشاہ} \ur{نے} \ur{ایک} \ur{نیا} \ur{حکم} \ur{دیا۔}
\end{quote}

\begin{quote}
\emph{bādshāh ne ek nayā hukm diyā.}
\end{quote}

\begin{quote}
The king gave a new order.
\end{quote}

Masculine.

\begin{cousinweb}
\textbf{\ur{حکم}} (\emph{hukm}) is usually traced, through Persian, to Arabic \emph{ḥukm}, \textquotedblleft{}a judgement, a command\textquotedblright{}, a cousin of \textbf{\ur{حکومت}}.
\end{cousinweb}

\begin{grammarlens}[title={What you've built}]
One more word for telling what happened.
\end{grammarlens}

\subsection*{Your turn}
\begin{itemize}
\raggedright
  \item \emph{Say it:} \emph{hukm}
  \item \emph{Say it:} \emph{hukm}, once more
  \item \emph{Say it:} \emph{bādshāh ne ek nayā hukm diyā.} — The king gave a new order.
  \item \emph{Recall:} say the Urdu for a peace settlement, then the Urdu for the subjects (of a ruler), then say \emph{hukm} again
\end{itemize}

\begin{tcolorbox}[breakable,colback=teal!4,colframe=teal!35!black,title={Before you move on}]
What does \ur{حکم} mean? (\textquotedblleft{}An order; a command\textquotedblright{}.) Say it once more.
\end{tcolorbox}
\section[\texorpdfstring{\ur{غلام} (ghulām)}{غلام (ghulām)}]{\ur{غلام} (ghulām) — a slave}

\label{lesson:UR-C312-ghulam}



\begin{quote}
Before the new one: say the Urdu for the subjects (of a ruler), then the Urdu for an order.
\end{quote}

\subsection*{You'll want to know: \ur{غلام}}
\textbf{\ur{غلام}} — \emph{ghulām} — \textquotedblleft{}a slave\textquotedblright{}.

\begin{quote}
\ur{پہلے} \ur{یہاں} \ur{غلام} \ur{کام} \ur{کرتے} \ur{تھے۔}
\end{quote}

\begin{quote}
\emph{pehle yahāṅ ghulām kām karte the.}
\end{quote}

\begin{quote}
Slaves used to work here long ago.
\end{quote}

Masculine.

\begin{cousinweb}
\textbf{\ur{غلام}} (\emph{ghulām}) is usually traced, through Persian, to Arabic \emph{ghulām}, \textquotedblleft{}a youth, a servant\textquotedblright{}; in Urdu it is one who is owned. \textbf{\ur{غلامی}} is slavery.
\end{cousinweb}

\begin{grammarlens}[title={What you've built}]
One more word for telling what happened.
\end{grammarlens}

\subsection*{Your turn}
\begin{itemize}
\raggedright
  \item \emph{Say it:} \emph{ghulām}
  \item \emph{Say it:} \emph{ghulām}, once more
  \item \emph{Say it:} \emph{pehle yahāṅ ghulām kām karte the.} — Slaves used to work here long ago.
  \item \emph{Recall:} say the Urdu for the subjects (of a ruler), then the Urdu for an order, then say \emph{ghulām} again
\end{itemize}

\begin{tcolorbox}[breakable,colback=teal!4,colframe=teal!35!black,title={Before you move on}]
What does \ur{غلام} mean? (\textquotedblleft{}A slave\textquotedblright{}.) Say it once more.
\end{tcolorbox}
\section[\texorpdfstring{\ur{بغاوت} (baghāvat)}{بغاوت (baghāvat)}]{\ur{بغاوت} (baghāvat) — a rebellion; a revolt}

\label{lesson:UR-C312-baghavat}



\begin{quote}
Before the new one: say the Urdu for an order, then the Urdu for a slave.
\end{quote}

\subsection*{You'll want to know: \ur{بغاوت}}
\textbf{\ur{بغاوت}} — \emph{baghāvat} — \textquotedblleft{}a rebellion; a revolt\textquotedblright{}.

\begin{quote}
\ur{شہر} \ur{میں} \ur{بغاوت} \ur{ہوئی۔}
\end{quote}

\begin{quote}
\emph{shahr meṅ baghāvat huī.}
\end{quote}

\begin{quote}
There was a rebellion in the city.
\end{quote}

Feminine.

\begin{cousinweb}
\textbf{\ur{بغاوت}} (\emph{baghāvat}) is usually traced, through Persian, to Arabic \emph{baghy}, \textquotedblleft{}a revolt\textquotedblright{}. It is feminine: \textbf{\ur{بغاوت} \ur{ہوئی}}.
\end{cousinweb}

\begin{grammarlens}[title={What you've built}]
One more word for telling what happened.
\end{grammarlens}

\subsection*{Your turn}
\begin{itemize}
\raggedright
  \item \emph{Say it:} \emph{baghāvat}
  \item \emph{Say it:} \emph{baghāvat}, once more
  \item \emph{Say it:} \emph{shahr meṅ baghāvat huī.} — There was a rebellion in the city.
  \item \emph{Recall:} say the Urdu for an order, then the Urdu for a slave, then say \emph{baghāvat} again
\end{itemize}

\begin{tcolorbox}[breakable,colback=teal!4,colframe=teal!35!black,title={Before you move on}]
What does \ur{بغاوت} mean? (\textquotedblleft{}A rebellion; a revolt\textquotedblright{}.) Say it once more.
\end{tcolorbox}
\section[\texorpdfstring{\ur{للکارنا} (lalkārnā)}{للکارنا (lalkārnā)}]{\ur{للکارنا} (lalkārnā) — to challenge; to call out (an enemy)}

\label{lesson:UR-C312-lalkarna}



\begin{quote}
Before the new one: say the Urdu for a slave, then the Urdu for a rebellion.
\end{quote}

\subsection*{You'll want to know: \ur{للکارنا}}
\textbf{\ur{للکارنا}} — \emph{lalkārnā} — \textquotedblleft{}to challenge; to call out (an enemy)\textquotedblright{}.

\begin{quote}
\ur{صبح} \ur{ایک} \ur{سپاہی} \ur{نے} \ur{بادشاہ} \ur{کو} \ur{للکارا۔}
\end{quote}

\begin{quote}
\emph{subah ek sipāhī ne bādshāh ko lalkārā.}
\end{quote}

\begin{quote}
In the morning a soldier challenged the king.
\end{quote}

Past, with \ur{نے} (he or she …): \ur{اس} \ur{نے} \ur{شیر} \ur{کو} \ur{للکارا}, \emph{us ne sher ko lalkārā}.

\begin{cousinweb}
\textbf{\ur{للکارنا}} (\emph{lalkārnā}) is built on \emph{lalkār}, \textquotedblleft{}a challenge, a defiance\textquotedblright{}, which CDIAL lists for Hindi under a form it reconstructs (CDIAL 10973). The one challenged takes \textbf{\ur{کو}}.
\end{cousinweb}

\begin{grammarlens}[title={What you've built}]
One more word for telling what happened.
\end{grammarlens}

\subsection*{Your turn}
\begin{itemize}
\raggedright
  \item \emph{Say it:} \emph{lalkārnā}
  \item \emph{Say it:} \emph{lalkārnā}, once more
  \item \emph{Say it:} \emph{subah ek sipāhī ne bādshāh ko lalkārā.} — In the morning a soldier challenged the king.
  \item \emph{Recall:} say the Urdu for a slave, then the Urdu for a rebellion, then say \emph{lalkārnā} again
\end{itemize}

\begin{tcolorbox}[breakable,colback=teal!4,colframe=teal!35!black,title={Before you move on}]
What does \ur{للکارنا} mean? (\textquotedblleft{}To challenge; to call out (an enemy)\textquotedblright{}.) Say it once more.
\end{tcolorbox}
